\documentclass[11pt]{article}
\usepackage[T1]{fontenc}
\usepackage{lmodern,amsmath,amssymb,amsthm}
\usepackage[margin=1in]{geometry}
\usepackage[hidelinks]{hyperref}
\newtheorem{theorem}{Theorem}
\newtheorem{proposition}[theorem]{Proposition}
\newcommand{\F}{\mathbb F}
\newcommand{\Q}{\mathrm{Q}}
\newcommand{\Zp}{\mathbb Z_p}
\newcommand{\Sel}{\operatorname{Sel}}
\newcommand{\unr}{\mathrm{ur}}
\title{A Selmer Dimension Cut at Resonant Primes for Split Cartan Normalizers}
\author{\small\href{mailto:admin@cosmosgeek.org}{admin@cosmosgeek.org}\\
\small\href{mailto:b321014@iiit-bh.ac.in}{b321014@iiit-bh.ac.in}}
\date{Working draft for specialist review}
\begin{document}\maketitle
\begin{abstract}
We treat a fixed-depth auxiliary-prime problem for an elliptic
Galois representation with image the normalizer of a split Cartan
modulo $p^m$, $p\ge7$. At a prime $q$ extending the swapping
element of this normalizer one has $q\equiv-1\bmod p^m$.
The adjoint tame space has two root lines. Relaxing to only one
of them allows the primal and dual residual Selmer dimensions
to be lowered by one simultaneously, under a detectable
off-diagonal projection hypothesis. The negative choice of
special Frobenius eigenvalues also preserves the original
two-dimensional representation's unramified local finite line
through the stated precision. We give the hypotheses for iteration
and the local calculation for a resulting special lift.
\end{abstract}

\section{The finite-depth problem}
Let $E/\Q$ be an elliptic curve, $p\ge7$, $k=\F_p$, and $m\ge1$.
Write $\rho_m$ for its action on $E[p^m]$ and $\bar\rho$ for
its reduction. Suppose there is an integral basis at this depth
in which
\[
 G_m=\operatorname{im}\rho_m
  =\bigl((\mathbb Z/p^m)^\times\times(\mathbb Z/p^m)^\times\bigr)
                         \rtimes\langle w\rangle,\qquad
 w=\begin{pmatrix}0&1\\1&0\end{pmatrix}.
\]
The diagonal subgroup defines a quadratic field $F$. Our hypothesis
specifies the actual division-field image; it can apply to a CM
elliptic curve at an ordinary split prime. It is stronger than
irreducibility modulo $p$. Let $\eta$ be the quadratic character
of $F/\Q$, and write $\epsilon$ for the mod-$p$ cyclotomic character.
Set $A=\operatorname{ad}^0\bar\rho$ with the conjugation action.
The trace pairing identifies its Tate dual $A^*$ with $A(1)$.
There is a direct sum of $G_{\Q}$-modules
\[
 A=k\eta\oplus M,\qquad
 k\eta=k\operatorname{diag}(1,-1),\qquad
 M=kE_{12}\oplus kE_{21}.
\]
Fix a finite set $\Sigma$ of places containing $\infty,p$ and all
places of ramification. At $\Sigma$ prescribe linear local
conditions $\mathcal L$ on $A$, and take unramified conditions
elsewhere. Let $\mathcal L^\perp$ be the exact local Tate-orthogonal
conditions on $A^*$. Assume
\[
 \dim_k\Sel_{\mathcal L}(A)
 =\dim_k\Sel_{\mathcal L^\perp}(A^*)=d,\qquad
 \Sel_{\mathcal L}(A)\cap H^1(\Q,k\eta)=0.
 \tag{1}
\]
The last condition refers to the indicated direct summand. One
sufficient criterion is that every pure $\eta$-Selmer class restrict
to an everywhere unramified class over $F$ and $p\nmid h_F$. Indeed,
restriction across the quadratic extension is injective for such
classes and class field theory identifies their images with a
subspace of $\operatorname{Hom}(\operatorname{Cl}(F),k)$.
To apply this criterion one must check the old local conditions,
including the one at $p$.

\begin{theorem}\label{cut}
Under (1), if $d>0$ there are infinitely many good primes
$q\notin\Sigma$ with $\rho_m(\operatorname{Frob}_q)$ conjugate to
$w$, at each of which one can replace the unramified adjoint
condition by the upper tame root line in the $(+1,-1)$ eigenbasis
of $w$. The resulting primal Selmer group is a codimension-one
subspace of the old group, and the resulting dual group also has
dimension $d-1$. There is a succession of $d$ distinct such
primes after which both groups vanish. They may avoid any
prescribed finite set. One may also impose splitting in a fixed
quadratic field $K$ ramified at a good prime outside $\Sigma$,
provided that prime is kept outside all added ramification sets.
\end{theorem}

The use of auxiliary-prime local conditions to vary Selmer groups
is classical in deformation theory; see \cite{CM,FKP}. The point at
$q\equiv-1\bmod p$ is that full local $H^1$ is three-dimensional.
The two-dimensional partial relaxation and the full depth-$m$
swapping class are both part of the construction.

\section{Simultaneous Frobenius detection}
In the diagonal Cartan the two off-diagonal lines have distinct
characters $s/t,t/s$, and $w$ exchanges them. Thus $M$ is an
irreducible $G_1$-module. Every scalar matrix fixes $M$, and the
scalar $zI$ acts on $A^*$ by $z^2$. Choosing $z^2\ne1$ shows
that no constituent of $A^*$ is isomorphic to $M$.

We need
\[
                 H^1(G_m,M)=H^1(G_m,A^*)=0.
\]
For $A^*$ the central scalar above gives the assertion by the
central-element cocycle identity. For $M$, the order of $G_1$ is
prime to $p$. The kernel of $G_m\to G_1$ is diagonal; conjugation
by diagonal matrices is trivial there. Its maximal elementary
abelian quotient therefore has no equivariant map to $M$, which has
no diagonal invariants. Inflation--restriction proves the assertion
also for $M$.

Fix $f\ne0$ in the primal group and $g\ne0$ in the dual group. By
(1) the projection $f_M$ is nonzero. Their restrictions to
$\ker\rho_m$ are additive homomorphisms; the restrictions of $f_M$
and $g$ are nonzero by the preceding vanishings. Their joint image
in $M\oplus A^*$ is a $G_1$-submodule. The projections are $M$ and
a nonzero submodule $C$ of $A^*$. This image equals $M\oplus C$:
the two projections of a subdirect product have a common quotient
after dividing by their kernels in the two summands, and here such
a nonzero quotient is excluded by the scalar action. In particular
one can alter independently the two translation coordinates of a
lift of $w$ in their joint finite cocycle extension, while leaving
its entire depth-$m$ matrix fixed.

In the eigenbasis $(1,1),(1,-1)$ of $w$, write $H',e',f'$ for the
trace-zero diagonal, upper and lower matrices. In these coordinates
the components in the original basis are
\[
 \operatorname{diag}(1,-1)=e'+f',\quad
 E_{12}+E_{21}=H',\quad E_{12}-E_{21}=-e'+f'.
\]
Frobenius on $A$ acts as $+1$ on $H'$ and $-1$ on the other two
vectors. Its cyclotomic twist reverses those signs. The unramified
value of an $A$-cocycle at this Frobenius is its projection onto
$kH'$. The diagonal $\eta$-component contributes zero, and the
arbitrary $M$-translation lets us make the $f$-value nonzero.
For a dual unramified value $B$, the linear functional needed
below is $\operatorname{tr}(e'B)$. It is nonzero on each simple
constituent of $A^*=k\eta(1)\oplus M(1)$: on the invariant line
in the first constituent use $e'+f'$, and in the second use
$-e'+f'$. The two constituents have different dimensions and are
irreducible. Hence $C$ can vary the $g$-value to make this functional
nonzero, while keeping the chosen $f$-value.

Chebotarev supplies infinitely many primes in the resulting conjugacy
class. The eigenlines in the condition are labelled by their distinct
eigenvalues; conjugating the data preserves the stated nonvanishing.
At a good prime the elliptic determinant is cyclotomic, so
$q\equiv-1\bmod p^m$. The prime is inert in $F$. If the extra field
$K$ is used, the joint finite division and cocycle field is
unramified at the prescribed good ramified place of $K$. Its
intersection with that quadratic field is therefore $\Q$,
so the independent splitting tag can be appended to Chebotarev.

\section{The local condition and the cut}
For finite unramified coefficients at $q\ne p$, restriction to tame
inertia identifies the quotient by unramified $H^1$ with the
$q$-eigenspace of Frobenius acting on the coefficient module.
This follows from inflation--restriction and the tame relation
$\varphi\tau\varphi^{-1}=\tau^q$. One may choose coordinates
using an arithmetic Frobenius lift. At the selected $q$ the relevant
spaces in these coordinates are
\[
 \begin{split}
 H^1(\Q_q,A)&= kH'_{\unr}\oplus k e'_{\rm tame}\oplus k f'_{\rm tame},\\
 H^1(\Q_q,A^*)&= k e'_{\unr}\oplus k f'_{\unr}\oplus kH'_{\rm tame}.
 \end{split}
\]
Let $U=kH'_{\unr}$ and $N=k e'_{\rm tame}$, the latter represented
by classes with Frobenius coordinate zero and this upper root on
inertia. It is the tangent to the chosen special condition
$\varphi\mapsto\operatorname{diag}(-q,-1)$,
$\tau\mapsto1+c e'$. Indeed the displayed diagonal conjugates $e'$
to $q e'$, so the tame relation holds for every $c$ and any lift
of $c$. This specifies one of the two possible root directions.
Local Tate duality pairs $e'_{\rm tame}$ with an unramified dual
class by the trace against $e'$ up to a fixed nonzero scalar:
pairing the tame generator with an unramified Frobenius generator
is nonzero in the local invariant map, and the coefficient pairing
is trace.

First relax $U$ to $U+N$ only. Pair any partially relaxed global
class with the \emph{old} dual class $g$. Local pairings at every
other place vanish. At $q$ the unramified component pairs to zero
and the pairing with $g$ is nonzero on the one-dimensional $N$.
The reciprocity sum forces the new class's component along $N$
to vanish. Thus the partially relaxed Selmer group already equals
the old group. Replacing $U$ by $N$ inside this relaxation now
keeps precisely those old classes with zero localization at $q$.
As $f$ has nonzero localization, that is a single independent
linear equation. The global duality dimension formula preserves
the difference of primal and dual dimensions: the only changed
local condition has the same dimension, one. This proves both
dimension bounds. The new primal group lies in the old one, so the
pure-diagonal exclusion persists. Induction proves
Theorem~\ref{cut} without assuming the mixed local condition
respects the original global decomposition $k\eta\oplus M$.

\section{The elliptic finite line at a negative root}
The local condition on the two-dimensional lattice is a separate
calculation. Let $T$ be a free $\Zp$-lattice with the local special
action above, $V=T\otimes\Q_p$, $D=V/T$. In a fixed basis
suppose Frobenius is $\operatorname{diag}(-q,-1)$ and inertia
acts by $\tau\mapsto 1+cE_{12}$ on a pro-$p$ tame generator
with $c\ne0$. Let $j\ge1$ satisfy $p^j\mid(q+1)$ and $p^j\mid c$.
Then $T/p^j$ is unramified, and the propagated finite condition in
$H^1(\Q_q,T/p^j)$ from $D$ is its full unramified line.

Indeed $V^I$ is spanned by the first basis vector and
$\varphi-1$ acts there by $-(q+1)\ne0$, so
$H^1_f(\Q_q,V)=H^1_{\unr}(\Q_q,V)=0$. Its image in $H^1(\Q_q,D)$
is zero. Under the identification $T/p^j\simeq D[p^j]$, the
preimage of that condition is the connecting image of
$D^{G_{\Q_q}}/p^jD^{G_{\Q_q}}$.
Frobenius acts with second eigenvalue $-1$, so invariants of $D$
lie on the first coordinate and form
$(q+1)^{-1}\Zp/\Zp$ there. A preimage under multiplication by $p^j$
of the generator $1/(q+1)$ may be taken on the same line. Its
connecting cocycle vanishes on inertia and has Frobenius value
$-1/p^j$ on that coordinate. This generates the unramified
cohomology, which is $(T/p^j)/(\varphi-1)(T/p^j)$ and has one
free coordinate. Thus if a global special lift retains the given
representation through level $p^j$, the old elliptic local line
agrees with this propagated one. Existence of a global
characteristic-zero modular lift with a specified complete set of
old local conditions is an additional deformation problem.

\begin{thebibliography}{9}
\bibitem{CM} F.~Calegari and B.~Mazur,
\emph{Nearly ordinary Galois deformations over arbitrary number fields},
J. Inst. Math. Jussieu \textbf{8} (2009), 99--177.
\href{https://www.math.uchicago.edu/~fcale/papers/ordinary.pdf}{Authors' paper}.
\bibitem{FKP} N.~Fakhruddin, C.~Khare and S.~Patrikis,
\emph{Relative deformation theory, relative Selmer groups, and lifting
irreducible Galois representations},
\href{https://par.nsf.gov/servlets/purl/10324335}{authors' paper}.
\bibitem{record} BSD research record, 16 November 2026,
checkpoint 46, \emph{CM-inert negative primes: balanced rigidity and
a primitive ordinary source}, sections 1--7 and independent review.
\end{thebibliography}
\end{document}
